\documentclass[11pt]{article}
\usepackage[T1]{fontenc}
\usepackage[utf8]{inputenc}
\usepackage{lmodern}
\usepackage[margin=1in]{geometry}
\usepackage{amsmath,amssymb,mathtools}
\usepackage[hidelinks]{hyperref}
\usepackage{xcolor}

\setlength{\parindent}{0pt}
\setlength{\parskip}{0.65em}
\definecolor{statusblue}{RGB}{24,73,120}
\newcommand{\Z}{\mathbb Z}
\newcommand{\Ent}{\operatorname{Ent}}

\title{The cyclic-group hypercontractivity problem is now complete\\
\large Literature answer to the open scope in arXiv:2512.03489}
\author{Run \texttt{fa\_banach\_001} --- lane 10}
\date{12 August 2026}

\begin{document}
\maketitle

\textbf{Status.} \textcolor{statusblue}{\texttt{literature\_already\_answered}.}
The cited papers give a full characterization for every cyclic order.
Yao's May 2026 continuation proves the sharp result for every \(n\ge4\), and
Xie--Zhang independently prove the same sharp log--Sobolev inequality.  A
February 2026 paper of Cao--Fan--Han--Qiu--Wang gives the different exact
answer for the exceptional group \(\Z_3\).  Together with the classical
two-point case, these papers determine the optimal time for every \(n\ge2\).

\section*{Question in the source}

For \(\psi_n(k)=\min(k,n-k)\), let the Poisson-like semigroup on \(\Z_n\) be
\[
 P_t\!\left(\sum_{k=0}^{n-1}a_k\chi_k\right)
 =\sum_{k=0}^{n-1}e^{-t\psi_n(k)}a_k\chi_k.
\]
For \(1<p\le q<\infty\), the problem is to determine the least \(t_{p,q}(n)\)
such that
\[
 \|P_tf\|_q\le\|f\|_p\qquad(t\ge t_{p,q}(n)).             \tag{1}
\]
The source \cite{Yao25} proves
\[
 t_{p,q}(n)=\frac12\log\frac{q-1}{p-1}
\]
for \(n=2^k\) and \(n=3\cdot2^k\), \(k\ge1\), via the sharp
log--Sobolev inequality
\[
 \Ent_{\mu_n}(f^2)\le 2\langle f,A_{\psi_n}f\rangle.    \tag{2}
\]
Its abstract then says that the general case for arbitrary \(n\) remains
open.

\section*{Explicit completion for every \(n\ge4\)}

Yao's continuation \cite[Theorems 1.1 and 1.2]{Yao26} states, for every
integer \(n\ge4\),
\[
 \boxed{
 \|P_tf\|_q\le\|f\|_p
 \quad\Longleftrightarrow\quad
 t\ge\frac12\log\frac{q-1}{p-1}}
 \qquad(1<p\le q<\infty),                               \tag{3}
\]
and proves (2) with the optimal constant \(2\).  This is explicitly presented
as the complete solution and cites arXiv:2512.03489 as the prior dyadic-tower
result.

The proof does not merely add more towers.  It reduces the \(n\)-LSI to the
nonnegativity of an explicit polynomial of degree at most three in
\(\lfloor n/2\rfloor\) variables, then treats all odd \(n\ge5\) and all even
\(n\ge4\).  Thus every cyclic order omitted by the source is covered.

There is also an independent confirmation.  Xie and Zhang
\cite[Theorem 1 and Corollary 2]{XZ} prove for every \(n\ge4\) the same sharp
inequality (2), and consequently exactly (3), using a different
cubic-majorant/Fourier-block method.  Their introduction explicitly says
that Yao obtained the result first and that their proof was carried out
independently.

\section*{The exceptional order \(n=3\)}

Order three is genuinely different: the threshold in (3) is generally
false.  To state the exact answer, write \(T_r\) for the multiplier
\(r^{\psi_3(k)}\).  Since \(P_t=T_{e^{-t}}\), if \(r_{p,q}(\Z_3)\) is the
largest \(r\) for which \(\|T_rf\|_q\le\|f\|_p\), then
\[
 t_{p,q}(3)=-\log r_{p,q}(\Z_3).                         \tag{4}
\]
Cao--Fan--Han--Qiu--Wang \cite[Theorem 1.2]{CFHQW} prove that, for
\(1<p<q<\infty\),
\[
 \boxed{r_{p,q}(\Z_3)=\frac{(1+2x)(1-y)}{(1+2y)(1-x)},} \tag{5}
\]
where \((x,y)\) is the unique solution in \((0,1)^2\) of
\[
 \frac1{1+2x}\left(\frac{1+2x^p}{3}\right)^{1/p}
 =\frac1{1+2y}\left(\frac{1+2y^q}{3}\right)^{1/q},      \tag{6}
\]
\[
 \frac{(1-x)(1-x^{p-1})}{1+2x^p}
 =\frac{(1-y)(1-y^{q-1})}{1+2y^q}.                      \tag{7}
\]
Equations (4)--(7) are an exact finite system and simultaneously identify a
nonconstant extremizer.  When \(p=q\), of course \(t_{p,p}(3)=0\).

\section*{Full classification and scope}

Combining the cited results gives
\[
 \boxed{
 t_{p,q}(n)=
 \begin{cases}
 \frac12\log\dfrac{q-1}{p-1},&n=2\text{ or }n\ge4,\\[6pt]
 -\log r_{p,q}(\Z_3),&n=3,
 \end{cases}}
 \qquad 1<p<q<\infty,                                  \tag{8}
\]
with \(r_{p,q}(\Z_3)\) given uniquely by (5)--(7).  This fully closes the
source's arbitrary-order hypercontractivity problem.  Formula (3) also gives
the full constant-\(2\) LSI answer requested in the source for all \(n\ge4\);
\(n=3\) is exceptional rather than a missing instance of that claim.

The cheap run indexes were searched for all four arXiv identifiers, titles,
cyclic hypercontractivity, and word-length LSI phrases; no prior packet was
found.  Current web searches on 12 August 2026 located the explicit completion,
the independent proof, and the exceptional-order characterization.

\begin{thebibliography}{9}
\bibitem{Yao25}
G. Yao, \emph{Optimal Hypercontractivity and Log--Sobolev inequalities on
Cyclic Groups \(\Z_{m\cdot2^k}\)}, arXiv:2512.03489 (2025).

\bibitem{Yao26}
G. Yao, \emph{A Complete Solution of Optimal Hypercontractivity and
Log--Sobolev inequalities on Cyclic Groups \(\Z_n\) for \(n\ge4\)},
\href{https://arxiv.org/abs/2605.30867}{arXiv:2605.30867} (2026).

\bibitem{XZ}
X. Xie and H. Zhang, \emph{Sharp Log--Sobolev Inequalities for Finite Cyclic
Groups with Word-Length},
\href{https://arxiv.org/abs/2606.02847}{arXiv:2606.02847} (2026).

\bibitem{CFHQW}
J. Cao, S. Fan, Y. Han, Y. Qiu, and Z. Wang,
\emph{The Optimal Hypercontractive Constants for \(\Z_3\) and Biased
Bernoulli Random Variables},
\href{https://arxiv.org/abs/2602.17248}{arXiv:2602.17248} (2026).
\end{thebibliography}

\end{document}
